\documentclass[12pt,a4paper]{article}

% Paquetes esenciales para codificación e idioma español
\usepackage[utf8]{inputenc}
\usepackage[T1]{fontenc}
\usepackage[spanish,es-tabla]{babel}

% Paquetes requeridos por tu código
\usepackage{booktabs}
\usepackage{longtable}
\usepackage{xcolor}
\usepackage{hyperref}

% Definiciones de comandos y colores personalizados para evitar errores
\definecolor{alerta}{RGB}{220, 53, 69}
\newcommand{\pendiente}{\textbf{--}}

\begin{document}
	
	% Sección padre requerida para que las \subsection tengan coherencia en la numeración
	\section{Arquitectura y Tecnologías}
	
	% ============================================================
	% SECCION 5.1 - Historia MVC, Front Controller y comparativa de
	% frameworks backend. Pegar dentro de \section{...} correspondiente
	% del informe. Requiere paquetes ya presentes en main.tex:
	% booktabs, longtable, hyperref, xcolor.
	% ============================================================
	
	\subsection{Historia del patrón MVC y el Front Controller}
	\label{sec:mvc-front-controller}
	
	El patrón \textbf{Modelo-Vista-Controlador (MVC)} fue formulado por Trygve
	Reenskaug en 1979 durante su estancia en Xerox PARC, como una forma de
	separar la representación interna de los datos (Modelo) de su presentación
	al usuario (Vista) y de la lógica que coordina ambas (Controlador) dentro
	del lenguaje Smalltalk-80. La idea original nació para interfaces gráficas
	de escritorio, no para la web: cada cambio en el Modelo notificaba a la
	Vista mediante el patrón Observador, y el Controlador traducía las
	acciones del usuario en cambios sobre el Modelo.
	
	A finales de los años 90 y principios de los 2000, frameworks web como
	Struts (Java) adaptaron MVC al ciclo petición/respuesta HTTP, que es
	fundamentalmente distinto al modelo de eventos de una GUI de escritorio:
	no hay un objeto Vista persistente que escuche al Modelo, sino que cada
	petición HTTP dispara una ejecución completa que termina generando una
	respuesta y muere. Esta adaptación introdujo el patrón
	\textbf{Front Controller} (Alur, Crupi y Malks, \emph{Core J2EE Patterns},
	2001): en lugar de que cada recurso (página, endpoint) maneje su propia
	entrada, un único punto de entrada centralizado recibe \emph{todas} las
	peticiones HTTP entrantes, resuelve routing, aplica lógica transversal
	(autenticación, logging, manejo de errores) y delega el trabajo específico
	al controlador correspondiente.
	
	Los frameworks modernos usados en el stack de este proyecto y en las
	alternativas evaluadas implementan el Front Controller de forma muy
	similar, aunque con nombres distintos:
	
	\begin{itemize}
		\item \textbf{Spring Boot 3}: el \texttt{DispatcherServlet} es el Front
		Controller; recibe toda petición HTTP, la enruta según las anotaciones
		\texttt{@RequestMapping}/\texttt{@GetMapping} de los controladores, y
		delega en \texttt{HandlerAdapter}.
		\item \textbf{Laravel 11}: \texttt{public/index.php} es el único punto
		de entrada; el \texttt{Kernel} HTTP procesa la petición a través de
		middleware y la enruta según \texttt{routes/web.php} o
		\texttt{routes/api.php}.
		\item \textbf{Symfony 7}: análogo a Laravel, \texttt{public/index.php}
		instancia el \texttt{Kernel}, que resuelve la ruta vía el componente
		\texttt{Routing} y ejecuta el controlador asociado.
		\item \textbf{ASP.NET Core 8}: el \emph{middleware pipeline} configurado
		en \texttt{Program.cs} actúa como Front Controller; el middleware de
		\texttt{Routing} despacha hacia el controlador o \emph{minimal API}
		correspondiente.
		\item \textbf{Django 5}: el manejador WSGI/ASGI central resuelve la URL
		contra \texttt{urls.py} y delega en la vista (\emph{view function} o
		\emph{class-based view}) correspondiente.
	\end{itemize}
	
	\subsubsection{Comparativa de frameworks backend}
	
	La Tabla~\ref{tab:comparativa-frameworks} resume las cinco opciones
	evaluadas para el backend de este proyecto. Las columnas de throughput
	provienen de la prueba \emph{Fortunes} (la más representativa de un CRUD
	real: ejercita ORM, conectividad a base de datos, ordenamiento y
	renderizado) de la Ronda 23 de TechEmpower Framework Benchmarks, publicada
	en febrero de 2025 y —hasta donde se pudo verificar al momento de escribir
	este informe— la última ronda completa antes de que el proyecto fuera
	discontinuado en marzo de 2026.
	
	\begin{table}[h!]
		\centering
		\small
		\begin{tabular}{@{}lllrl@{}}
			\toprule
			\textbf{Framework} & \textbf{Lenguaje} & \textbf{Versión evaluada} & \textbf{Fortunes (req/s)} & \textbf{Front Controller} \\
			\midrule
			ASP.NET Core & C\# & 8 & $\sim$609\,966 & Middleware pipeline \\
			Spring Boot  & Java & 3 & $\sim$243\,639 & \texttt{DispatcherServlet} \\
			Symfony      & PHP  & 7 & N/D\footnotemark & \texttt{Kernel} + \texttt{Routing} \\
			Django       & Python & 5 & $\sim$32\,651  & Resolver WSGI/ASGI \\
			Laravel      & PHP  & 11 & $\sim$16\,800   & \texttt{Kernel} HTTP \\
			\bottomrule
		\end{tabular}
		\caption{Comparativa de frameworks backend --- prueba Fortunes, TechEmpower Round 23 (feb. 2025).}
		\label{tab:comparativa-frameworks}
	\end{table}

	Es importante leer esta tabla con cautela académica: los benchmarks de
	TechEmpower miden un límite superior de throughput bajo condiciones de
	laboratorio (hardware dedicado, sin lógica de negocio real, implementaciones
	optimizadas por la comunidad para el propio benchmark), no el rendimiento
	que tendría cada framework en un proyecto con la complejidad de SGED. Aun
	así, la tendencia relativa —plataformas compiladas/JIT con E/S asíncrona
	(ASP.NET Core, Spring Boot) por delante de frameworks interpretados
	orientados a productividad (Laravel, Django)— es consistente con la
	literatura y con benchmarks independientes.
	
	\subsubsection{Justificación de ADR-001}
	
	La decisión registrada en \textbf{ADR-001} (ver
	\texttt{docs/adr/ADR-001-tecnologia.md}) elige Java 21 con Spring Boot 3.2.5
	sobre Laravel 11 y ASP.NET Core 8, las otras dos opciones válidas indicadas
	por el docente. Los criterios de esa decisión no fueron puramente de
	rendimiento bruto —de hecho, la Tabla~\ref{tab:comparativa-frameworks}
	muestra que ASP.NET Core supera a Spring Boot en throughput teórico—, sino
	un balance entre:
	
	\begin{enumerate}
		\item \textbf{Coherencia de stack tipado}: Angular (TypeScript) y Spring
		Boot (Java) comparten filosofía de tipado estático, lo que reduce
		errores de integración frente a un backend débilmente tipado.
		\item \textbf{Seguridad integrada}: Spring Security ofrece soporte JWT y
		control de acceso por roles sin librerías de terceros adicionales.
		\item \textbf{Costo de herramientas}: IntelliJ IDEA Community Edition es
		gratuito, a diferencia de Visual Studio completo para .NET.
		\item \textbf{Curva de aprendizaje aceptable} dentro de las 10 semanas
		disponibles, mitigada por Spring Initializr y la amplia documentación
		del ecosistema.
	\end{enumerate}
	
	En otras palabras, el throughput medido por TechEmpower es un dato de
	contexto relevante para justificar por qué Spring Boot es una opción
	razonable en términos de rendimiento (no la más lenta del grupo evaluado),
	pero no fue el criterio decisivo de ADR-001.
	
	% ============================================================
	% SUBSECCION 5.4 - Apache Bench. PLANTILLA: completar la tabla y el
	% parrafo de analisis con los datos reales generados por
	% scripts/bench-ab.sh una vez el backend este corriendo localmente.
	% Los valores salen de docs/mediciones/perf/ab/ab-resumen-<fecha>.md
	% ============================================================
	
%	\subsection{Carga concurrente con Apache Bench}
%	\label{sec:apache-bench}
%	
%	Como complemento a la comparativa teórica de la sección anterior, se midió
%	el comportamiento del backend propio bajo carga concurrente usando
%	\textbf{Apache Bench (ab)}, ejecutado mediante
%	\texttt{scripts/bench-ab.sh}. El script automatiza: autenticación contra
%	\texttt{POST /api/auth/login}, extracción de la cookie JWT
%	(\texttt{sged\_access}), y el disparo de \texttt{ab -n 1000 -c 50} contra
%	un endpoint autenticado (\texttt{GET /api/estudiantes}).
%	
%	\begin{table}[h!]
%		\centering
%		\begin{tabular}{@{}lr@{}}
%			\toprule
%			\textbf{Métrica} & \textbf{Valor} \\
%			\midrule
%			Peticiones totales & \pendiente{} \\
%			Concurrencia        & \pendiente{} \\
%			Peticiones fallidas / \% error & \pendiente{} \\
%			Peticiones por segundo & \pendiente{} \\
%			P50 (ms) & \pendiente{} \\
%			P95 (ms) & \pendiente{} \\
%			P99 (ms) & \pendiente{} \\
%			\bottomrule
%		\end{tabular}
%		\caption{Resultados de \texttt{ab -n 1000 -c 50} contra \texttt{GET /api/estudiantes}. \textbf{PENDIENTE: reemplazar con datos reales de \texttt{docs/mediciones/perf/ab/}.}}
%		\label{tab:apache-bench}
%	\end{table}
	
\end{document}